\documentclass{article}
\usepackage{tcee-repaso}
\tcetema[Economía del bienestar (I)]{3.A.22}{Economía del bienestar (I). Los teoremas fundamentales del bienestar. Óptimo económico y «second-best».}
\tcefecha{21/10/2026}
\begin{document}
\cabecerarepaso
\begin{columnas}

% ==========================================================================
\bloque{Esquema del cante}
\esquemacante{%
\apartadocante{Introducción: lo mejor y lo bueno}{2,5}
\apartadocante{I. Óptimo económico. Criterio de \autor{Pareto}}{9,5}
\subcante{1. Definición y juicios de valor}{1}
\subcante{2. Condiciones: producción, consumo, global; GFPU}{6,5}
\subcante{3. Valoración del criterio}{2}
\apartadocante{II. Teoremas fundamentales}{10}
\subcante{1. 1TFEB: EGC paso a paso, prueba, evidencia}{5,5}
\subcante{2. 2TFEB: hiperplano separador, límites}{4,5}
\apartadocante{III. La teoría del \emph{second best}}{7}
\subcante{1. Idea y resultado (\autor{Lipsey}-\autor{Lancaster})}{3}
\subcante{2. Aplicaciones}{1,5}
\subcante{3. Información privada: mecanismos}{1,5}
\subcante{4. Valoración}{1}
\apartadocante{Conclusión}{1}}
\emph{Hilo}: el refrán «lo mejor es enemigo de lo bueno»: \hyperref[rep:I]{I} define lo mejor; \hyperref[rep:II]{II}, cuándo lo alcanza el mercado; \hyperref[rep:III]{III}, qué es lo bueno si lo mejor no está al alcance. \emph{Pizarra}: las 3 condiciones recuadradas hasta el final de \hyperref[rep:II]{II}; debajo, los 4 pasos del EGC, cada uno junto a su condición. Gráficos imprescindibles: GFPU, caja de los teoremas y los 2 del \emph{second best}.

% ==========================================================================
\bloque[2,5]{Introducción}
\begin{itemize}
\item \textbf{Enganche}: «Lo mejor es enemigo de lo bueno» (\autor{Voltaire}, 1772): la economía del bienestar lo convierte en teorema. CBAM, régimen definitivo desde el 1/1/2026: lo mejor sería un precio mundial único del CO\textsubscript{2} = daño marginal; como no existe, la UE grava el carbono de las importaciones en frontera: \emph{segundo óptimo}.
\item \textbf{Relevancia}: \emph{teórica}: \concepto{economía del bienestar} = parte normativa de la micro: valora los estados de la economía (asignación + reparto) y busca normas de política; \autor{Pigou} (1920): el origen de la economía es el «entusiasmo social». \emph{De política}: toda reforma, impuesto o regulación lleva juicios de bienestar.
\item \textbf{Contextualización}: \autor{Bentham} (1789, utilitarismo) $\to$ marginalistas (UMg cardinal; igualar ↑ bienestar) $\to$ \autor{Pareto} (1906): ordinal $\to$ nueva economía del bienestar (años 30-40): compensación (\autor{Kaldor}, \autor{Hicks}, 1939) y FBS (\autor{Bergson}, 1938; \autor{Samuelson}, 1947); condiciones (\autor{Lerner}, \autor{Lange}) $\to$ teoremas (\autor{Arrow}, \autor{Debreu}) $\to$ \emph{second best} $\to$ información (\autor{Mirrlees}, \autor{Hurwicz}). Teórica: completa el EG \vertema{3.A.21}: aquel, si existe, es único y estable; este, si es \emph{bueno}.
\item \textbf{Problemática}: ¿qué condiciones para que nadie mejore sin que otro empeore? ¿Las alcanza el mercado sin coordinador? ¿Puede llevar a cualquier óptimo? Si una no puede cumplirse, ¿conviene cumplir las demás?
\item \textbf{Estructura}: \hyperref[rep:I]{I}, óptimo (criterio, condiciones 2×2×2×2 y GFPU, valoración); \hyperref[rep:II]{II}, teoremas (1TFEB y 2TFEB); \hyperref[rep:III]{III}, \emph{second best} (idea, aplicaciones, información, valoración).
\end{itemize}

% ==========================================================================
\bloque[9,5]{I. Óptimo económico. Criterio de \autor{Pareto}\label{rep:I}}
\begin{ideaclave}
Un óptimo exige igualar 3 tipos de relaciones marginales (producción, consumo y entre ambas), sin precios ni mercados; pero hay infinitos, toda la GFPU, y la eficiencia no basta para elegir.
\end{ideaclave}

\sub{I.1. Definición y juicios de valor\label{rep:I.1}}
\begin{itemize}
\item \emph{Origen}: \autor{Pareto} (Lausana, sucesor de \autor{Walras} desde 1893; \emph{Manuale}, 1906): valorar el EG sin medir ni comparar utilidades. Versión moderna, con conjuntos y sin derivadas: \autor{Arrow} (1951), \autor{Debreu} (1959).
\item \concepto{Asignación}: reparto $\vect{x}_i$ + planes $\vect{y}_j$; \emph{factible} si $\sum_i\vect{x}_i\leq\sum_i\vect{\omega}_i+\sum_j\vect{y}_j$, $\vect{y}_j\in\mathcal{Y}_j$.
\item \concepto{\autor{Pareto} superior}: nadie peor y al menos uno mejor (nadie se opondría). \concepto{Óptimo de \autor{Pareto}} (eficiente): no existe ninguna factible \autor{Pareto} superior: no se mejora a nadie sin empeorar a otro. Superior $\neq$ óptimo (puede haber una tercera). Optimalidad \emph{fuerte} (aquí) implica la \emph{débil} (ninguna mejora a todos).
\item \emph{Ordinal}: solo si cada uno prefiere, no cuánto ni entre personas \vertema{3.A.8}.
\item \emph{4 juicios de valor}: i) \concepto{individualismo} (solo cuentan las utilidades); ii) \concepto{no paternalismo} (cada uno, mejor juez de lo suyo); iii) \concepto{benevolencia} (si uno mejora y nadie empeora, la sociedad mejora); iv) \concepto{independencia del proceso} (mercado, plan o imposición: da igual).
\end{itemize}

\sub{I.2. Condiciones de optimalidad paretiana\label{rep:I.2}}
\begin{itemize}
\item 3 condiciones: eficiencia en la \emph{producción}, en el \emph{consumo} (intercambio) y \emph{global}; las 2 primeras, necesarias pero no suficientes.
\item \emph{Modelo 2×2×2×2} (\autor{Bator}, 1957), estático, misma notación que \vertema{3.A.21}: 2 bienes $X$, $Y$; 2 factores $L$, $K$ con $\overline{L}$, $\overline{K}$; 2 consumidores $A$, $B$ (preferencias completas, transitivas, continuas, estrict. monótonas y convexas; $U_i$ estrict. cuasicóncavas y diferenciables); 2 empresas, $F_X$, $F_Y$ con isocuantas convexas y conjuntos de producción convexos; sin externalidades. \concepto{Planificador benevolente} que lo sabe todo: solución \emph{tecnocrática}.
\item \emph{No se supone}: UMg decreciente (sin sentido con utilidad ordinal) ni \autor{Inada}. Basta la convexidad: CPO suficientes y FPP cóncava.
\item \emph{Producción}: $\max F_X$ s.a. $F_Y(\overline{L}-L_X,\overline{K}-K_X)\geq\overline{Y}$:
\begin{formula}
\centering $\mathrm{RMST}^X_{LK}\equiv\dfrac{\mathrm{PMg}^X_L}{\mathrm{PMg}^X_K}=\dfrac{\mathrm{PMg}^Y_L}{\mathrm{PMg}^Y_K}\equiv\mathrm{RMST}^Y_{LK}$
\end{formula}
\item[] tangencia de isocuantas: \concepto{curva de contrato} de la producción ($O_X$ a $O_Y$) $\to$ \concepto{FPP}. $\lambda=\mathrm{PMg}^X_L/\mathrm{PMg}^Y_L$: inversa de la pendiente de la FPP. No es en general la diagonal; vértices también eficientes.
\item \emph{FPP}: i) pendiente negativa; su valor absoluto, la \concepto{relación marginal de transformación} (coste de oportunidad de $X$ en $Y$):
\begin{formula}
\centering $\mathrm{RMT}_{XY}\equiv-\dfrac{dY}{dX}\Big|_{\text{FPP}}=\dfrac{\mathrm{PMg}^Y_L}{\mathrm{PMg}^X_L}=\dfrac{\mathrm{CMg}_X}{\mathrm{CMg}_Y}$
\end{formula}
\item[] ii) forma: rend. decrecientes $\Rightarrow$ estrict. cóncava; RCE: cóncava si difieren las intensidades, recta si coinciden; rend. crecientes: puede ser convexa; iii) todos sus puntos, eficientes en la producción.
\item \emph{Consumo}: $\max U_A$ s.a. $U_B(\overline{X}-X_A,\overline{Y}-Y_A)\geq\overline{U}_B$:
\begin{formula}
\centering $\mathrm{RMS}^A_{XY}\equiv\dfrac{\partial U_A/\partial X_A}{\partial U_A/\partial Y_A}=\dfrac{\partial U_B/\partial X_B}{\partial U_B/\partial Y_B}\equiv\mathrm{RMS}^B_{XY}$
\end{formula}
\item[] si no, hay un intercambio que mejora a los 2. Tangencia de curvas de indiferencia: curva de contrato ($O_A$ a $O_B$) $\to$ \concepto{FPU} en $(U_A,U_B)$ (\autor{Allais}, 1943; \autor{Samuelson}, 1950). $\lambda=-dU_A/d\overline{U}_B$.
\item \emph{FPU}: i) decreciente; ii) \emph{forma indeterminada}: utilidad ordinal $\Rightarrow$ una transformación monótona cambia su curvatura (se dibuja cóncava por costumbre); iii) todos sus puntos, eficientes en el consumo; iv) una FPU por cada punto de la FPP.
\item \emph{Global}: el planificador maximiza $U_A$ con $U_B\geq\overline{U}_B$, tecnologías y dotaciones (multiplicadores $\lambda_U,\lambda_X,\lambda_Y,\lambda_L,\lambda_K$). CPO: $\partial U_A/\partial X_A=\lambda_X=\lambda_U\partial U_B/\partial X_B$\ldots; $\lambda_X\mathrm{PMg}^X_L=\lambda_L=\lambda_Y\mathrm{PMg}^Y_L$\ldots; dividiendo, las \emph{3 condiciones de \autor{Pareto}}:
\funcionclave{\begin{aligned}&\mathrm{RMST}^X_{LK}=\mathrm{RMST}^Y_{LK}=\tfrac{\lambda_L}{\lambda_K}\\ &\mathrm{RMS}^A_{XY}=\mathrm{RMS}^B_{XY}=\tfrac{\lambda_X}{\lambda_Y}=\mathrm{RMT}_{XY}\end{aligned}}
\item[] Si $\mathrm{RMS}>\mathrm{RMT}$, se valora $X$ en más de lo que cuesta: producir más $X$. Multiplicadores = \concepto{precios sombra}; si los precios fueran proporcionales a ellos, cada agente elegiría lo que el planificador le asigna (\autor{Lange}-\autor{Lerner}): la puerta de \hyperref[rep:II]{II}. Con cuasiconcavidad y concavidad, suficientes. Eficiente en producción y consumo puede no serlo globalmente: falta $\mathrm{RMS}=\mathrm{RMT}$.
\item \emph{Gráfico}: un punto $E$ de la FPP fija la caja del consumo $(X^*,Y^*)$; en su curva de contrato solo $E^*$ tiene $\mathrm{RMS}^A=\mathrm{RMS}^B=\mathrm{RMT}$; cada punto de la FPP ($E_1,E_2,E_3$) da una FPU con su $E^*$; su envolvente: la \concepto{gran frontera de posibilidades de utilidad} (GFPU).
\end{itemize}
\graficorepaso[0.98\linewidth]{gfpu}{Construcción de la GFPU: caja del consumo dentro de la FPP con $\mathrm{RMS}=\mathrm{RMT}$ (izq.) y envolvente de las FPU (dcha.)}
\begin{itemize}
\item \emph{Forma}: utilidad ordinal $\Rightarrow$ ni las FPU ni la GFPU \emph{tienen por qué ser cóncavas}: solo decreciente; puede tener tramos convexos (\emph{hondonadas}).
\item \emph{Implicación}: infinitos óptimos (toda la GFPU, 10 valores cada uno) y \autor{Pareto} no los ordena: optimalidad \emph{necesaria pero no suficiente} para maximizar el bienestar social.
\item \concepto{Función de bienestar social} $W(U_A,U_B)$ (\autor{Bergson}, 1938; \autor{Samuelson}, 1947), creciente; curvas de indiferencia social. \concepto{Punto de felicidad} (\emph{bliss point}) $c$: tangencia con la $W$ más alta; con hondonadas, máximo \emph{global}. Un punto ineficiente ($a$, en $W_2$) puede ser mejor que uno eficiente y desigual ($d$, en $W_1$); pero siempre hay uno en la GFPU que lo mejora. Formas de $W$: \vertema{3.A.24}.
\end{itemize}
\graficorepaso[0.55\linewidth]{fbs-gfpu}{FBS y GFPU: punto de felicidad $c$}

\sub{I.3. Valoración del criterio de \autor{Pareto}\label{rep:I.3}}
\begin{itemize}
\item \emph{Ventajas}: ordinal; pocos juicios de valor, aceptados (no desperdiciar); permite valorar el mercado (\hyperref[rep:II]{II}).
\item \emph{Críticas}: i) \emph{estático} (austriacos: eficiencia dinámica, \autor{Huerta de Soto}, 2004); ii) \emph{favorece el statu quo}: veto de cada uno (unanimidad); iii) \emph{independiente del proceso}; iv) \emph{ignora la distribución}: que uno lo tenga todo es óptimo (\autor{Sen}, 1970: lujo e inanición; asignaciones \emph{justas} = eficientes + libres de envidia, \autor{Foley}, 1967); v) \emph{incompleto} (la principal): no ordena los óptimos ni compara uno eficiente con uno ineficiente que no esté abajo a la izquierda (\autor{Wicksell}: confundir óptimo competitivo con óptimo social).
\item \emph{Respuestas} (años 30-40): \emph{neoparetiana} (LSE): \concepto{criterios de compensación}: ganadores \emph{podrían} compensar a perdedores (\autor{Kaldor}, \autor{Hicks}, 1939); potencial, pueden contradecirse (\autor{Scitovsky}, 1941); análisis coste-beneficio \vertema{3.A.9}. \emph{Neoutilitarista} (Harvard): FBS de \autor{Bergson}-\autor{Samuelson}, juicios explícitos \vertema{3.A.24}.
\end{itemize}

% ==========================================================================
\bloque[10]{II. Teoremas fundamentales\label{rep:II}}
\begin{ideaclave}
EGC y óptimo se corresponden en los 2 sentidos: todo equilibrio es un óptimo (1TFEB) y todo óptimo puede ser un equilibrio tras redistribuir (2TFEB); cada sentido exige supuestos distintos, y en ellos está la política económica.
\end{ideaclave}
\begin{itemize}
\item \concepto{Equilibrio walrasiano} \vertema{3.A.21}: $\vect{p}^*$ y $(\vect{x}^*_i,\vect{y}^*_j)$ con: i) cada empresa maximiza $\vect{p}^*\cdot\vect{y}_j$ en $\mathcal{Y}_j$; ii) cada consumidor maximiza $U$ con $\overline{W}_i=\vect{p}^*\cdot\vect{\omega}_i+\sum_j\theta_{ij}\vect{p}^*\cdot\vect{y}^*_j$; iii) se vacían todos los mercados. Con \concepto{transferencias de suma fija} $T_i$ ($\sum T_i=0$): \emph{equilibrio con transferencias}.
\item \emph{Origen}: \autor{Lerner} y \autor{Lange} (cálculo socialista); \autor{Hotelling} (1938), \autor{Allais} (1943); análisis convexo: \autor{Arrow} y \autor{Debreu}.
\end{itemize}

\sub{II.1. Primer teorema (1TFEB)\label{rep:II.1}}
\begin{itemize}
\item \emph{Enunciado} (\autor{Arrow}, 1951; \autor{Debreu}, 1951): con preferencias \emph{localmente no saciables}, todo equilibrio walrasiano es un óptimo de \autor{Pareto}. Formaliza la \emph{mano invisible} de \autor{Smith} (1776).
\item \emph{Supuestos} (misma economía de \autor{Bator} + mercados con $p_X,p_Y,w_L,w_K$ dados; $\overline{W}_A=w_L\overline{L}_A+w_K\overline{K}_A+\theta_{AX}B_X+\theta_{AY}B_Y$):
\end{itemize}
{\centering\small\renewcommand{\arraystretch}{1.05}\setlength{\tabcolsep}{2pt}%
\begin{tabular}{@{}>{\raggedright\arraybackslash}m{0.42\linewidth}>{\raggedright\arraybackslash}m{0.54\linewidth}@{}}
\toprule
\rowcolor{tceeazulpalido}\textbf{Supuesto} & \textbf{Si se levanta (fallo)}\\
\midrule
1. No saciedad local & renta sin gastar: puede no ser eficiente\\
2. Mercados completos & externalidades, bienes públicos \vertema{3.A.23}\\
3. Precio-aceptantes & poder de mercado \vertema{3.A.17}-\vertema{3.A.19}\\
4. Precios comunes & impuestos: exceso de gravamen \vertema{4.B.9}\\
5. Información perfecta & selección adversa, riesgo moral \vertema{3.A.13}\\
\bottomrule
\end{tabular}\par}
\begin{itemize}
\item \emph{No hace falta}: convexidad ni condición sobre rendimientos (aseguran que el equilibrio \emph{exista}); el teorema: \emph{si existe}, es eficiente. Derivadas y soluciones interiores, solo para el desarrollo.
\item \emph{Paso 1, consumidores} (cumplen la de consumo): $\max U_A$ s.a. $p_XX_A+p_YY_A\leq\overline{W}_A$:
\begin{formula}
\centering $\mathrm{RMS}^A_{XY}=\dfrac{p_X}{p_Y}=\mathrm{RMS}^B_{XY}$
\end{formula}
\item[] mismos precios $\Rightarrow$ valoran igual sin conocerse.
\item \emph{Paso 2, empresas en los factores} (la de producción): $\max p_jF_j-w_LL_j-w_KK_j$ $\Rightarrow$ $p_j\mathrm{PMg}^j_L=w_L$, $p_j\mathrm{PMg}^j_K=w_K$:
\begin{formula}
\centering $\mathrm{RMST}^X_{LK}=\dfrac{w_L}{w_K}=\mathrm{RMST}^Y_{LK}$
\end{formula}
\item \emph{Paso 3, empresas en los bienes} (con los consumidores, la global): $p_j=w_L/\mathrm{PMg}^j_L=\mathrm{CMg}_j$ \vertema{3.A.16}:
\begin{formula}
\centering $\dfrac{p_X}{p_Y}=\dfrac{\mathrm{CMg}_X}{\mathrm{CMg}_Y}=\mathrm{RMT}_{XY}$
\end{formula}
\item[] nadie sabe de tecnología ni de preferencias ajenas.
\item \emph{Paso 4, vaciado}: $X_A+X_B=F_X$, $Y_A+Y_B=F_Y$, $L_X+L_Y=\overline{L}$, $K_X+K_Y=\overline{K}$; ley de \autor{Walras}: 3 bastan, precios relativos (numerario). El vaciado convierte las 3 igualdades en una asignación: sin él, planes incompatibles.
\item \emph{Resultado}: las condiciones de \hyperref[rep:I.2]{I.2} con $p_X/p_Y=\lambda_X/\lambda_Y$ y $w_L/w_K=\lambda_L/\lambda_K$: \emph{los precios de mercado hacen de precios sombra}. Con cuasiconcavidad y concavidad, suficientes: 1TFEB en 2×2×2×2.
\item \emph{Gráfico}: desde $\vect{\omega}$, el equilibrio $E$ está en la curva de contrato: la recta de precios es tangente a las 2 curvas en el mismo punto; con producción, $p_X/p_Y$ = pendiente común de la FPP en $E$ y de las curvas en $E^*$.
\end{itemize}
\graficorepaso[0.68\linewidth]{caja-intercambio-teoremas}{1TFEB: $E$ en la curva de contrato; 2TFEB: de $\vect{\omega}$ a $\vect{\omega}'$ para alcanzar $E'$}
\begin{itemize}
\item \emph{Implicaciones}: i) \emph{economía de información}: el planificador necesita todo; a cada agente le bastan los precios y lo suyo (\autor{Hayek}, 1945); ii) \emph{papel del Estado}: corregir solo los fallos que rompen un supuesto (defensa del libre mercado); iii) llega a \emph{un} óptimo, no \emph{al} óptimo: el de las dotaciones, no tiene por qué ser el punto de felicidad.
\item \emph{Prueba general} (\autor{Arrow}-\autor{Debreu}; $I$, $J$, $n$; sin derivadas, vale en esquinas), por reducción al absurdo: si una factible es \autor{Pareto} superior, quien mejora no podía pagarla ($>\overline{W}_i$) y quien queda igual, por no saciedad local, tampoco por menos ($\geq\overline{W}_i$); sumando, y como las empresas maximizan el beneficio, su valor supera el de lo disponible: contradice la factibilidad. Vale con transferencias de suma fija.
\item \emph{Evidencia}: subastas dobles de \autor{V. Smith} (1962): con pocos y solo su valoración o coste, convergen al equilibrio competitivo (Nobel 2002); \autor{Gode}-\autor{Sunder} (1993): con robots al azar, eficiencia casi completa: lo hacen las reglas del mercado.
\item \emph{Valoración}: nada de equidad; supuestos exigentes: cada \concepto{fallo de mercado} rompe uno (2: \vertema{3.A.23}; 3: poder; 4: impuestos; 5: información; ni óptimo restringido, \autor{Greenwald}-\autor{Stiglitz}, 1986); presupone existencia, sin unicidad ni estabilidad; estático.
\end{itemize}

\sub{II.2. Segundo teorema (2TFEB)\label{rep:II.2}}
\begin{itemize}
\item \emph{Enunciado} (\autor{Arrow}, 1951; \autor{Debreu}, 1954): con preferencias y tecnologías \emph{convexas} (y preferencias continuas y localmente no saciables), todo óptimo de \autor{Pareto} puede alcanzarse como equilibrio competitivo tras una \emph{redistribución de suma fija}. Recíproco del 1TFEB: no hace falta sustituir el mercado.
\item \emph{Supuestos}: i) los del 1TFEB; ii) \emph{convexidad} de preferencias y conjuntos de producción (sin rend. crecientes): el nuevo y decisivo; iii) transferencias de suma fija $T_i$; iv) \emph{condición de renta}: cada uno puede comprar algo más barato que su cesta óptima.
\item \emph{Intuición}: en $E'$ las curvas $U_A'$ y $U_B'$ son tangentes y la recta común las separa: es la recta de precios; el Estado lleva $\vect{\omega}$ a cualquier $\vect{\omega}'$ de esa recta y el mercado hace el resto. En utilidades, la economía se mueve por la GFPU de $E$ al punto de felicidad $E'$.
\end{itemize}
\pareja{segundo-teorema-gfpu}{2TFEB en utilidades: de $E$ al punto de felicidad $E'$}{segundo-teorema-no-convexo}{Sin convexidad: la tangente no sostiene $E'$}
\begin{itemize}
\item \emph{Demostración} (4 pasos, MWG §16.D): i) $V=\sum_iV_i$ (lo preferido estrictamente) y $\mathcal{Y}=\sum_j\mathcal{Y}_j+\{\sum_i\vect{\omega}_i\}$, convexos; ii) disjuntos (si no, el punto no sería óptimo); iii) \concepto{teorema del hiperplano separador} (\autor{Minkowski}, 1911): $\vect{p}\neq\vect{0}$ y $r$ con $\vect{p}\cdot\vect{z}\geq r$ en $V$, $\leq r$ en $\mathcal{Y}$: ese $\vect{p}$ son los precios; iv) \concepto{cuasiequilibrio} + condición de renta $\Rightarrow$ equilibrio, con $T_i=\vect{p}\cdot\vect{x}^*_i-\vect{p}\cdot\vect{\omega}_i-\sum_j\theta_{ij}\vect{p}\cdot\vect{y}^*_j$, $\sum T_i=0$.
\item \emph{Sin convexidad}: lo preferido «envuelve» la recta de precios; con rend. crecientes, el óptimo no maximiza el beneficio a ningún precio \vertema{3.A.17}.
\item \emph{Implicaciones}: \emph{separación eficiencia-equidad} (distribución con suma fija, asignación al mercado); intervención \emph{ex ante} por equidad aunque no haya fallos (punto de felicidad de una FBS \vertema{3.A.24}); \emph{descentralización}: socialismo de mercado (\autor{Lange}, 1936; \autor{Lerner}, 1944).
\item \emph{Límites}: i) \emph{información}: la suma fija que redistribuye de verdad depende de lo inobservable (capacidad, esfuerzo); si depende de lo observable, distorsiona (\autor{Mirrlees}, 1971): vuelve el dilema; ii) \emph{convexidad}: rend. crecientes $\Rightarrow$ $p=\mathrm{CMg}$ con pérdidas; iii) \emph{cálculo}: fijar $\vect{\omega}'$ exige la información que el 1TFEB ahorraba; iv) \emph{multiplicidad}: otros equilibrios desde $\vect{\omega}'$.
\item \emph{Soluciones}: imposición óptima sobre bienes (\autor{Ramsey}, 1927) y renta (\autor{Mirrlees}, 1971) \vertema{4.B.11}; eficiencia en la producción (\autor{Diamond}-\autor{Mirrlees}, 1971); gravar renta y no bienes con ocio separable (\autor{Atkinson}-\autor{Stiglitz}, 1976); \emph{tagging} (\autor{Akerlof}, 1978) \vertema{4.B.3}. Sin suma fija: lógica del \emph{second best}.
\end{itemize}

% ==========================================================================
\bloque[7]{III. La teoría del \emph{second best}\label{rep:III}}
\begin{ideaclave}
Si una condición de \autor{Pareto} no puede cumplirse, cumplir las demás no garantiza acercarse al óptimo: lo bueno no es lo que más se parece a lo mejor. Con información privada, lo alcanzable es un óptimo restringido por los incentivos.
\end{ideaclave}
\begin{itemize}
\item \concepto{Primer óptimo} (\emph{first best}): cumple todas las condiciones; \concepto{segundo óptimo} (\emph{second best}): el mejor con una restricción adicional. Nunca «de primer o segundo orden».
\end{itemize}

\sub{III.1. Idea y resultado general\label{rep:III.1}}
\begin{itemize}
\item \emph{Visión tradicional}: cuantas más condiciones, más cerca; eliminar distorsiones una a una (\emph{piecemeal}). \autor{Meade} (1955), uniones aduaneras, lo intuye; \autor{Lipsey}-\autor{Lancaster} (1956) lo demuestran: no necesariamente.
\item \emph{Supuestos}: $n$ bienes, consumidor representativo $U(\vect{x})$, tecnología $T(\vect{x})=0$; primer óptimo: $U_k/U_n=T_k/T_n$ $\forall k$. \concepto{Restricción no trivial} (efectiva y no suprimible) sobre el bien 1:
\begin{formula}
\centering $\dfrac{U_1}{U_n}=\psi\,\dfrac{T_1}{T_n},\quad\psi\neq1$\quad (impuesto \emph{ad valorem}: $\psi=1+t$)
\end{formula}
\item[] origen: institucional (impuesto, arancel, regulación), tecnológico (costes fijos del monopolio natural), de mercado (poder, externalidad), informativo.
\item \emph{Desarrollo}: $\max U$ s.a. $T=0$ ($\lambda_1$) y la restricción ($\lambda_2$); con $\Psi_k$ = derivada de $\frac{U_1}{U_n}-\psi\frac{T_1}{T_n}$:
\begin{formula}
\centering $\dfrac{U_k}{U_n}=\dfrac{T_k+\frac{\lambda_2}{\lambda_1}\Psi_k}{T_n+\frac{\lambda_2}{\lambda_1}\Psi_n}$
\end{formula}
\item[] $\Psi_k$: derivadas cruzadas de segundo orden, signo desconocido a priori.
\item \concepto{Teorema general del segundo óptimo} (\autor{Lipsey}-\autor{Lancaster}, 1956): con una restricción que impide cumplir una condición, las demás, aunque alcanzables, dejan en general de ser deseables; hay que apartarse también de ellas, en una dirección que depende de sustituciones y complementariedades.
\item \emph{Gráfico}: la restricción (recta) deja fuera $F$ (felicidad, en $W_3$); $P$, corte con la GFPU, cumple todas las condiciones posibles ($W_1$); $Q$, tangencia de la restricción con $W_2>W_1$, \emph{dentro} de la GFPU, es mejor: el segundo óptimo.
\end{itemize}
\pareja{second-best-gfpu}{\emph{Second best}: $Q$ (ineficiente) mejor que $P$ (eficiente accesible)}{montana-second-best}{Montaña: con el obstáculo $D$, mejor la cumbre $A$ que seguir hacia $V$}
\begin{itemize}
\item \emph{Montaña}: no se llega a lo más alto andando siempre cuesta arriba.
\item \emph{Implicaciones}: i) contar condiciones no informa; ii) quitar una distorsión puede ↑, ↓ o no cambiar el bienestar; iii) \emph{añadir} una puede mejorarlo si compensa otra (cine gravado y plataformas sustitutivas); iv) la política por partes no tiene garantía general.
\item \emph{Extensiones}: \emph{separabilidad} (\autor{Davis}-\autor{Whinston}, 1965): sector distorsionado separable ($\Psi_k=\Psi_n=0$) $\Rightarrow$ \autor{Pareto} en el resto: rescata el \emph{piecemeal}; \emph{eficiencia en la producción} (\autor{Diamond}-\autor{Mirrlees}, 1971): no gravar bienes intermedios; \emph{corregir en el origen} (\autor{Bhagwati}-\autor{Ramaswami}, 1963): un arancel ante distorsión interna es solo segundo óptimo y puede empeorar.
\end{itemize}

\sub{III.2. Aplicaciones\label{rep:III.2}}
\begin{itemize}
\item \emph{Uniones aduaneras} \vertema{3.B.10}: quitan aranceles entre socios y los mantienen fuera: creación frente a desviación de comercio (\autor{Viner}, 1950; \autor{Lipsey}, 1957); \autor{Meade}: expansión.
\item \emph{Monopolio contaminante} \vertema{3.A.17}\vertema{3.A.23}: produce poco (poder) y mucho (externalidad); el pigouviano reduce aún más $q$. Con $\mathrm{IMg}=\mathrm{CMgP}+t$ y $p(q^*)=\mathrm{CMgP}+\mathrm{CMgE}$:
\begin{formula}
\centering $t^*=\mathrm{CMgE}(q^*)-\dfrac{p(q^*)}{|\varepsilon_{q,p}|}$
\end{formula}
\item[] menor que el pigouviano; subvención si pesa más el poder (\autor{Buchanan}, 1969; \autor{Barnett}, 1980). Romper el monopolio sin gravar la contaminación puede ↓ bienestar \vertema{3.A.20}.
\item \emph{Industria naciente} \vertema{3.B.8}: corregir el fallo de capitales o de apropiación con crédito o subvención; el arancel, segundo óptimo (encarece el consumo).
\item \emph{Instituciones} \vertema{3.B.28}: \autor{Rodrik} (2008): instituciones de segundo óptimo, no copiar «mejores prácticas»; FMI, Hong Kong, sept.\ 1997: liberalizar capitales en plena crisis asiática \vertema{3.B.17}.
\item \emph{Imposición y regulación}: gravar más los complementos del ocio (\autor{Corlett}-\autor{Hague}, 1953) \vertema{4.B.11}; \autor{Ramsey}-\autor{Boiteux} y $p=\mathrm{CMe}$ del impugnable \vertema{3.A.17}\vertema{3.A.15}; variedad en competencia monopolística \vertema{3.A.18}.
\end{itemize}

\sub{III.3. Información privada: diseño de mecanismos\label{rep:III.3}}
\begin{itemize}
\item \concepto{Eficiencia restringida}: ningún planificador con las mismas restricciones de información y mercados mejora a todos; la restricción son los incentivos \vertema{3.A.13}. Mercados incompletos: abrir uno más puede empeorar a todos (\autor{Hart}, 1975); equilibrios no eficientes ni en sentido restringido (\autor{Greenwald}-\autor{Stiglitz}, 1986): la mano invisible es invisible porque no está (\autor{Stiglitz}, 2002).
\item \concepto{Diseño de mecanismos} (\autor{Hurwicz}, 1972): reglas para que, cada uno con su tipo $t_i$ y en su interés, se alcance el resultado. \emph{Principio de revelación}: basta mecanismos con $\mathrm{IC}_i$ y $\mathrm{PC}_i$ \vertema{3.A.14}.
\item \emph{VCG} (\autor{Vickrey}, 1961; \autor{Clarke}, 1971; \autor{Groves}, 1973): verdad dominante y eficiente (cada uno paga el daño que causa); utilidad cuasilineal; sin presupuesto equilibrado (\autor{Green}-\autor{Laffont}, 1977). Sin pagos: \autor{Gibbard}-\autor{Satterthwaite} \vertema{3.A.24}. \emph{Bayesianos} (\autor{d'Aspremont}-\autor{Gérard-Varet}, 1979; \autor{Arrow}, 1979): eficientes y equilibrados, no siempre voluntarios. \autor{Myerson}-\autor{Satterthwaite} (1983): en el intercambio bilateral, ninguno cumple las 4 (eficiente, veraz, equilibrado, voluntario).
\end{itemize}

\sub{III.4. Valoración\label{rep:III.4}}
\begin{itemize}
\item Sin reglas generales (signo de $\Psi_k$): dice qué no afirmar, no qué hacer; exige conocer todas las sustituciones; \emph{el mapa imposible} (\autor{J. Robinson}, 1962): un modelo con toda la realidad sería tan inútil como un mapa a escala 1:1.
\item \emph{Respuestas}: análisis parcial si las interdependencias son pequeñas; separabilidad; corregir en el origen; \autor{Lipsey} (2007): conocer el caso. No invalida los teoremas: enseña a usarlos.
\end{itemize}

% ==========================================================================
\bloque[1]{Conclusión}
\begin{itemize}
\item \textbf{Recapitulación}: óptimo = 3 igualdades marginales, infinitos (GFPU); 1TFEB: con mercados completos, precio-aceptación, precios comunes e información, el mercado alcanza uno mirando precios; 2TFEB: con convexidad, cualquiera tras suma fija; si una condición falla, cumplir las demás no garantiza nada (\emph{second best}).
\item \textbf{Relevancia}: vara de medir de la micro normativa: intervenir por eficiencia (falla el 1.º), por equidad (2.º), siempre con restricciones.
\item \textbf{Extensiones}: fallos de mercado \vertema{3.A.23}; FBS e imposibilidad de \autor{Arrow} (1951) \vertema{3.A.24}; imposición óptima \vertema{4.B.11}; mecanismos y subastas \vertema{3.A.14}.
\item \textbf{Opinión}: los teoremas ordenan las preguntas: qué supuesto falla, qué restricción no se quita, qué información hace falta.
\item \textbf{Idea final}: \autor{Pigou} (entusiasmo social) y \autor{Pareto} (atender a consecuencias morales y políticas): lo mejor es enemigo de lo bueno si se persigue a ciegas, pero sin saber cuál es lo mejor no se encuentra lo bueno.
\end{itemize}

% ==========================================================================
\bloque{Autores y fechas}
\hito{1776}{\autor{Smith}: la mano invisible}
\hito{1906}{\autor{Pareto}: \emph{Manuale}, criterio ordinal}
\hito{1920}{\autor{Pigou}: economía del bienestar}
\hito{1934}{\autor{Lerner} (y \autor{Lange}, 1942): condiciones}
\hito{1938}{\autor{Bergson}: FBS; \autor{Hotelling}: prueba}
\hito{1939}{\autor{Kaldor}, \autor{Hicks}: compensación}
\hito{1941}{\autor{Scitovsky}: paradoja de la reversión}
\hito{1943}{\autor{Allais}: FPU y prueba}
\hito{1947}{\autor{Samuelson}: FBS (\emph{Foundations})}
\hito{1951}{\autor{Arrow}, \autor{Debreu}: 1TFEB y 2TFEB (\autor{Debreu}, 1954)}
\hito{1955}{\autor{Meade}: uniones aduaneras}
\hito{1956}{\autor{Lipsey}-\autor{Lancaster}: \emph{second best}}
\hito{1957}{\autor{Bator}: modelo 2×2×2×2}
\hito{1962}{\autor{V. Smith}: subastas dobles}
\hito{1963}{\autor{Bhagwati}-\autor{Ramaswami}: corregir en el origen}
\hito{1965}{\autor{Davis}-\autor{Whinston}: separabilidad}
\hito{1971}{\autor{Mirrlees}; \autor{Diamond}-\autor{Mirrlees}}
\hito{1972}{\autor{Hurwicz}: mecanismos; Nobel a \autor{Arrow} y \autor{Hicks}}
\hito{1983}{\autor{Myerson}-\autor{Satterthwaite}}
\hito{1986}{\autor{Greenwald}-\autor{Stiglitz}}
\hito{2008}{\autor{Rodrik}: instituciones de segundo óptimo}

\bloque{Datos clave}
\begin{itemize}
\item CBAM: \dato{régimen definitivo desde el 1/1/2026}{UE, 2023}: acero, aluminio, cemento, fertilizantes, electricidad, hidrógeno.
\item Juego del ultimátum, metaanálisis de 75 resultados: \dato{oferta media 40\,\%, 16\,\% de rechazos}{Oosterbeek et al., 2004}.
\item FMI, Comité Provisional en Hong Kong: \dato{sept.\ 1997}{FMI, 1997}, con la crisis asiática en marcha desde julio.
\item Premios Nobel: \autor{Arrow} y \autor{Hicks} (EG y bienestar), \autor{Sen} (bienestar), \autor{V. Smith} (experimentos): \dato{1972, 1998, 2002}{Fundación Nobel}.
\end{itemize}

% ==========================================================================
\bloque{Preguntas del tribunal}
\begin{enumerate}
\item \textbf{¿Por qué no se supone UMg decreciente? ¿Forma de la FPU?} Con utilidad ordinal no tiene sentido; la FPU solo es decreciente, curvatura indeterminada.
\item \textbf{¿Qué agente cumple cada condición?} Consumidores la de consumo; empresas (factores) la de producción; consumidores + empresas (bienes, $p=\mathrm{CMg}$) la global; el vaciado las hace compatibles.
\item \textbf{¿Necesitan convexidad los teoremas?} El 1.º no (solo no saciedad local); el 2.º sí (hiperplano separador) y suma fija.
\item \textbf{¿Por qué no hay suma fija? ¿Qué se hace?} Depende de lo inobservable (\autor{Mirrlees}); imposición óptima, \emph{tagging}.
\item \textbf{¿Puede un impuesto ↑ el bienestar?} Sí: \emph{second best} (complementos del ocio, cine-plataformas; $t^*<$ pigouviano al monopolio).
\item \textbf{¿Segundo óptimo y óptimo restringido?} El 1.º, distorsión dada (\autor{Lipsey}-\autor{Lancaster}); el 2.º, restricción de incentivos e información.
\item \textbf{¿Qué criterio usa el ACB?} Compensación potencial (\autor{Kaldor}-\autor{Hicks}): no se paga y puede contradecirse (\autor{Scitovsky}).
\end{enumerate}

\begin{trampa}
\begin{itemize}
\item 1TFEB: \emph{no} necesita convexidad ni rendimientos (solo no saciedad local + mercados completos, precio-aceptación, precios comunes, información); la convexidad es para la \emph{existencia} y para el 2TFEB.
\item 2TFEB: convexidad + transferencias de \emph{suma fija} + condición de renta; redistribuye \emph{dotaciones}, no precios.
\item Las 2 primeras condiciones son necesarias, no suficientes: falta $\mathrm{RMS}=\mathrm{RMT}$.
\item $\mathrm{RMT}_{XY}=\mathrm{CMg}_X/\mathrm{CMg}_Y=\mathrm{PMg}^Y_L/\mathrm{PMg}^X_L$ (índices cruzados).
\item FPU y GFPU \emph{no tienen por qué ser cóncavas} (utilidad ordinal): solo decrecientes; la FPP sí es cóncava con convexidad.
\item El 1TFEB da \emph{un} óptimo, no el punto de felicidad; un punto ineficiente puede tener más bienestar que uno eficiente.
\item \emph{Second best} $\neq$ «condiciones de segundo orden»; cumplir más condiciones no es estar mejor.
\item Segundo óptimo (distorsión dada) $\neq$ óptimo restringido (información).
\item VCG: veraz y eficiente pero sin presupuesto equilibrado; \autor{Myerson}-\autor{Satterthwaite}: ninguno cumple las 4.
\end{itemize}
\end{trampa}

\end{columnas}
\end{document}
